\textbf{Phần 1. Trắc nghiệm lựa chọn}
\setcounterpageref{ex}{0}

\begin{ex} %[3D8N48]
Làm tròn số $43$ đến hàng chục ta được số
\choice
{\True $40$}
{$50$}
{$45$}
{$30$}
\loigiai{
Chữ số hàng đơn vị là 3 ($3 < 5$) nên làm tròn xuống $40$.\\
\textbf{Đáp án của câu hỏi là A}
}
\end{ex}

\begin{ex} %[3D8N48]
Làm tròn số $78$ đến hàng chục ta được số
\choice
{$70$}
{$75$}
{\True $80$}
{$90$}
\loigiai{
Chữ số hàng đơn vị là 8 ($8 \ge 5$) nên làm tròn lên $80$.\\
\textbf{Đáp án của câu hỏi là C}
}
\end{ex}

\begin{ex} %[3D8H48]
Làm tròn số $354$ đến hàng chục ta được số
\choice
{\True $350$}
{$360$}
{$400$}
{$300$}
\loigiai{
Hàng đơn vị là 4 ($< 5$) nên làm tròn thành $350$.\\
\textbf{Đáp án của câu hỏi là A}
}
\end{ex}

\begin{ex} %[3D8H48]
Làm tròn số $675$ đến hàng chục ta được số
\choice
{$670$}
{\True $680$}
{$700$}
{$600$}
\loigiai{
Hàng đơn vị là 5 ($\ge 5$) nên làm tròn lên $680$.\\
\textbf{Đáp án của câu hỏi là B}
}
\end{ex}

\begin{ex} %[3D8H48]
Làm tròn số $430$ đến hàng trăm ta được số
\choice
{$500$}
{$450$}
{\True $400$}
{$300$}
\loigiai{
Hàng chục là 3 ($< 5$) nên làm tròn thành $400$.\\
\textbf{Đáp án của câu hỏi là C}
}
\end{ex}

\begin{ex} %[3D8H48]
Làm tròn số $780$ đến hàng trăm ta được số
\choice
{$700$}
{\True $800$}
{$750$}
{$900$}
\loigiai{
Hàng chục là 8 ($\ge 5$) nên làm tròn lên $800$.\\
\textbf{Đáp án của câu hỏi là B}
}
\end{ex}

\begin{ex} %[3D8V48]
Làm tròn số $3\,462$ đến hàng chục ta được số
\choice
{$3470$}
{$3500$}
{$3400$}
{\True $3460$}
\loigiai{
Hàng đơn vị là 2 ($< 5$) nên được $3460$.\\
\textbf{Đáp án của câu hỏi là D}
}
\end{ex}

\begin{ex} %[3D8V48]
Làm tròn số $5\,681$ đến hàng trăm ta được số
\choice
{$5600$}
{\True $5700$}
{$5680$}
{$6000$}
\loigiai{
Hàng chục là 8 ($\ge 5$) nên làm tròn lên $5700$.\\
\textbf{Đáp án của câu hỏi là B}
}
\end{ex}

\begin{ex} %[3D8V48]
Một đoạn đường dài $1\,248\text{ m}$. Làm tròn độ dài đoạn đường đến hàng trăm ta được
\choice
{$1200\text{ m}$}
{$1300\text{ m}$}
{\True $1200\text{ m}$ (vì chữ số hàng chục là $4 < 5$)}
{$1250\text{ m}$}
\loigiai{
Hàng chục là 4 ($< 5$) nên làm tròn thành $1200\text{ m}$.\\
\textbf{Đáp án của câu hỏi là C}
}
\end{ex}

\begin{ex} %[3D8V48]
Một trường tiểu học có $856$ học sinh. Làm tròn số học sinh đến hàng trăm ta được khoảng
\choice
{$800$ học sinh}
{$860$ học sinh}
{$850$ học sinh}
{\True $900$ học sinh}
\loigiai{
Hàng chục là 5 ($\ge 5$) nên làm tròn lên $900$ học sinh.\\
\textbf{Đáp án của câu hỏi là D}
}
\end{ex}

\begin{ex} %[3D8V48]
Có bao nhiêu số tự nhiên có 2 chữ số khi làm tròn đến hàng chục thì được số 40?
\choice
{\True $10$ số (từ $35$ đến $44$)}
{$9$ số}
{$11$ số}
{$8$ số}
\loigiai{
Các số từ 35 đến 44 gồm có: $44 - 35 + 1 = 10$ số.\\
\textbf{Đáp án của câu hỏi là A}
}
\end{ex}

\begin{ex} %[3D8C48]
Tìm số lớn nhất có 3 chữ số mà khi làm tròn đến hàng trăm thì được số 600.
\choice
{$650$}
{$640$}
{$599$}
{\True $649$}
\loigiai{
Để làm tròn đến hàng trăm được 600 thì chữ số hàng chục phải $< 5$. Số lớn nhất là $649$.\\
\textbf{Đáp án của câu hỏi là D}
}
\end{ex}

\textbf{Phần 2. Trắc nghiệm đúng sai}
\setcounterpageref{ex}{0}

\begin{ex} %[3D8H48]
Xét tính đúng sai của việc làm tròn số:
\choiceTF
{\True Làm tròn $56$ đến hàng chục được $60$.}
{\True Làm tròn $350$ đến hàng trăm được $400$.}
{Làm tròn $749$ đến hàng trăm được $800$.}
{\True Làm tròn $749$ đến hàng trăm được $700$ (vì hàng chục là $4 < 5$).}
\loigiai{
\begin{itemize}
    \item $56 \rightarrow 60$: Đúng.
    \item $350 \rightarrow 400$: Đúng.
    \item Hàng chục là 4 nên làm tròn thành 700. Ý 3 Sai, Ý 4 Đúng.
\end{itemize}
\begin{itemize}
    \item a) \textbf{Đúng}.
    \item b) \textbf{Đúng}.
    \item c) \textbf{Sai}.
    \item d) \textbf{Đúng}.
\end{itemize}
}
\end{ex}

\begin{ex} %[3D8V48]
Cho số $4\,573$. Xét tính đúng sai:
\choiceTF
{\True Làm tròn đến hàng chục được $4\,570$.}
{\True Làm tròn đến hàng trăm được $4\,600$.}
{Làm tròn đến hàng chục được $4\,580$.}
{\True Chữ số hàng đơn vị là $3 < 5$ nên giữ nguyên hàng chục.}
\loigiai{
\begin{itemize}
    \item Hàng chục: $4570$. Đúng.
    \item Hàng trăm: $4600$ (vì hàng chục là $7 \ge 5$). Đúng.
    \item Được 4570 (không phải 4580). Sai.
    \item Giải thích đúng: Đúng.
\end{itemize}
\begin{itemize}
    \item a) \textbf{Đúng}.
    \item b) \textbf{Đúng}.
    \item c) \textbf{Sai}.
    \item d) \textbf{Đúng}.
\end{itemize}
}
\end{ex}

\textbf{Phần 3. Trả lời ngắn}
\setcounterpageref{ex}{0}

\begin{ex} %[3D8H48]
Làm tròn số $85$ đến hàng chục được số nào?
\shortans{90}
\loigiai{
\textbf{Đáp số:} $90$\\
$90$.
}
\end{ex}

\begin{ex} %[3D8H48]
Làm tròn số $430$ đến hàng trăm được số nào?
\shortans{400}
\loigiai{
\textbf{Đáp số:} $400$\\
$400$.
}
\end{ex}

\begin{ex} %[3D8V48]
Làm tròn số $2\,368$ đến hàng chục được số nào?
\shortans{2370}
\loigiai{
\textbf{Đáp số:} $2370$\\
$2370$.
}
\end{ex}

\begin{ex} %[3D8V48]
Làm tròn số $7\,852$ đến hàng trăm được số nào?
\shortans{7900}
\loigiai{
\textbf{Đáp số:} $7900$\\
$7900$.
}
\end{ex}

\begin{ex} %[3D8V48]
Tìm số bé nhất có 3 chữ số mà khi làm tròn đến hàng trăm thì được số 500.
\shortans{450}
\loigiai{
\textbf{Đáp số:} $450$\\
$450$.
}
\end{ex}

\begin{ex} %[3D8C48]
Có bao nhiêu số tròn chục khi làm tròn đến hàng trăm được số 800?
\shortans{10}
\loigiai{
\textbf{Đáp số:} $10$\\
Các số tròn chục từ 750 đến 840 gồm có 10 số.
}
\end{ex}
